\documentclass{article}
\usepackage{amsmath, amssymb, amsthm}
\usepackage{fontspec}
\setmainfont{Times New Roman}

\newtheorem{theorem}{Твърдение}
\newtheorem{definition}{Дефиниция}
\newtheorem{corollary}{Следствие}

\title{Пораждащи функции и основни класове случайни величини}
\author{}
\date{}

\begin{document}
\maketitle

\section{Пораждащи функции}

В тази глава ще разгледаме пораждащите функции, които представляват полезна трансформация на случайните величини. Подобно на архитектурна скица, която дава различен поглед върху една сграда в зависимост от ъгъла, пораждащата функция предоставя мощен аналитичен инструмент за изследване на свойствата на дискретните разпределения.

\begin{definition}
Нека $X : \Omega \rightarrow \mathbb{N} \cup \{0\}$ е целочислена случайна величина (т.е. възможните ѝ стойности са $0, 1, 2, \dots$). Тогава под \textbf{пораждаща функция} (probability generating function) на $X$ разбираме функцията $g_X(s)$, дефинирана като:
\[ g_X(s) := \mathbb{E}s^X = \sum_{n=0}^{\infty} s^n \mathbb{P}(X=n), \quad \text{за } s \in [-1, 1]. \]
\end{definition}

С помощта на пораждащата функция могат лесно да се пресмятат математическото очакване и дисперсията на случайната величина, както и да се възстановят нейните вероятности.

\begin{theorem}
Нека $X$ е целочислена случайна величина с пораждаща функция $g_X(s)$. Тогава:
\begin{enumerate}
    \item[а)] $\mathbb{E}X = g_X'(1)$
    \item[б)] $\mathbb{D}X = g_X''(1) + g_X'(1) - (g_X'(1))^2$
    \item[в)] $k! \mathbb{P}(X=k) = g_X^{(k)}(0)$, за $k \ge 0$.
\end{enumerate}
\end{theorem}

\begin{proof}
а) По дефиниция $g_X(s) = \sum_{n=0}^{\infty} s^n \mathbb{P}(X=n)$. Диференцирайки по $s$, получаваме:
\[ g_X'(s) = \frac{d}{ds} \sum_{n=0}^{\infty} s^n \mathbb{P}(X=n) = \sum_{n=1}^{\infty} n s^{n-1} \mathbb{P}(X=n) = \mathbb{E}[X s^{X-1}]. \]
Замествайки $s=1$, намираме $g_X'(1) = \sum_{n=1}^{\infty} n \mathbb{P}(X=n) = \mathbb{E}X$.

б) Аналогично, втората производна е:
\[ g_X''(s) = \mathbb{E}[X(X-1)s^{X-2}]. \]
При $s=1$ получаваме $g_X''(1) = \mathbb{E}[X(X-1)] = \mathbb{E}X^2 - \mathbb{E}X$. 
От формулата за дисперсия $\mathbb{D}X = \mathbb{E}X^2 - (\mathbb{E}X)^2$, можем да изразим $\mathbb{E}X^2 = g_X''(1) + \mathbb{E}X = g_X''(1) + g_X'(1)$. Така:
\[ \mathbb{D}X = g_X''(1) + g_X'(1) - (g_X'(1))^2. \]

в) За да намерим вероятността $X$ да приеме стойност $k$, диференцираме $g_X(s)$ $k$ пъти. Всички членове със степен по-малка от $k$ се анулират, а членът със степен $k$ става $k! \mathbb{P}(X=k)$. При заместване с $s=0$, всички останали членове (съдържащи $s$) се нулират, и остава точно $k! \mathbb{P}(X=k)$.
\end{proof}

\begin{corollary}
Нека $X$ и $Y$ са целочислени случайни величини. Тогава $X \stackrel{d}{=} Y$ (равни по разпределение) тогава и само тогава, когато $g_X \equiv g_Y$.
\end{corollary}

Важно свойство на пораждащите функции се проявява при суми от независими случайни величини. 

\begin{definition}
Дискретните случайни величини $X_1, \dots, X_n$ (или $(X_j)_{j=1}^\infty$) имат \textbf{еднакво разпределение}, ако $X_i \stackrel{d}{=} X_j$ за всеки $i, j$.
\end{definition}

\begin{theorem}
Нека $X_1, X_2, \dots, X_n$ са независими в съвкупност целочислени случайни величини. Ако $Y = \sum_{j=1}^n X_j$, то:
\[ g_Y(s) = \prod_{j=1}^n g_{X_j}(s). \]
Ако в допълнение $X_j$ са еднакво разпределени, то $g_Y(s) = (g_{X_1}(s))^n$.
\end{theorem}
\begin{proof}
Ще го покажем за $n=2$. Нека $Y = X_1 + X_2$. Тогава:
\[ g_Y(s) = \mathbb{E}s^{X_1+X_2} = \sum_{j=0}^{\infty} \sum_{i=0}^{\infty} s^{i+j} \mathbb{P}(X_1=j, X_2=i). \]
Поради независимостта на $X_1$ и $X_2$, съвместната вероятност се разпада: $\mathbb{P}(X_1=j, X_2=i) = \mathbb{P}(X_1=j)\mathbb{P}(X_2=i)$.
\[ g_Y(s) = \sum_{j=0}^{\infty} s^j \mathbb{P}(X_1=j) \sum_{i=0}^{\infty} s^i \mathbb{P}(X_2=i) = g_{X_1}(s) g_{X_2}(s). \]
\end{proof}

\section{Схема на Бернули}

Схемата на Бернули е базов вероятностен модел, описващ поредица от независими експерименти с два възможни изхода (успех или неуспех) и постоянна вероятностна структура.
Разглеждаме редица $(X_j)_{j=1}^{\infty}$, които са независими в съвкупност и еднакво разпределени:
\begin{align*}
\mathbb{P}(X_j=1) &= p \quad (\text{вероятност за успех}) \\
\mathbb{P}(X_j=0) &= 1-p = q \quad (\text{вероятност за неуспех})
\end{align*}

\subsection{Разпределение на Бернули}
Случайната величина $X_1$ се нарича разпределена по Бернули, $X_1 \sim \text{Ber}(p)$.
\begin{center}
\begin{tabular}{c|c|c}
$X_1$ & 0 & 1 \\ \hline
$\mathbb{P}$ & $q$ & $p$
\end{tabular}
\end{center}
Пораждащата ѝ функция е:
\[ g_{X_1}(s) = q \cdot s^0 + p \cdot s^1 = q + ps. \]
Математическото очакване и дисперсията са:
\[ \mathbb{E}X_1 = p, \quad \mathbb{D}X_1 = \mathbb{E}X_1^2 - (\mathbb{E}X_1)^2 = p - p^2 = pq. \]

\section{Биномно разпределение}

Биномното разпределение описва броя на успехите в първите $n$ експеримента от схема на Бернули. Записваме $X \sim \text{Bi}(n, p)$.
Случайната величина $X$ може да се представи като сума от независими бернулиеви величини: $X = \sum_{j=1}^n X_j$.

\begin{theorem}
Нека $X \sim \text{Bi}(n, p)$. Тогава:
\begin{enumerate}
    \item[а)] $g_X(s) = (q+ps)^n$
    \item[б)] $\mathbb{E}X = np$
    \item[в)] $\mathbb{D}X = npq$
    \item[г)] $\mathbb{P}(X=k) = \binom{n}{k} p^k q^{n-k}$, за $0 \le k \le n$.
\end{enumerate}
\end{theorem}
\begin{proof}
а) Тъй като $X = \sum_{j=1}^n X_j$ и $X_j$ са независими и еднакво разпределени, то $g_X(s) = (g_{X_1}(s))^n = (q+ps)^n$.

б) $\mathbb{E}X = g_X'(1) = np(q+ps)^{n-1} \big|_{s=1} = np(p+q)^{n-1} = np$.

г) $\mathbb{P}(X=k)$ може да се намери чрез $k$-тата производна в нулата:
\[ k! \mathbb{P}(X=k) = g_X^{(k)}(0) = \left. \frac{d^k}{ds^k} (q+ps)^n \right|_{s=0} = n(n-1)\dots(n-k+1) p^k q^{n-k}. \]
Следователно $\mathbb{P}(X=k) = \frac{n(n-1)\dots(n-k+1)}{k!} p^k q^{n-k} = \binom{n}{k} p^k q^{n-k}$.
\end{proof}

\section{Геометрично разпределение}

Геометричното разпределение моделира броя на неуспехите до настъпване на първия успех в схема на Бернули. Записваме $X \sim \text{Ge}(p)$. Формално:
\[ X = \min\left\{n \ge 1 : \sum_{j=1}^n X_j = 1\right\} - 1. \]

\begin{theorem}
Нека $X \sim \text{Ge}(p)$. Тогава:
\begin{enumerate}
    \item[а)] $\mathbb{P}(X=k) = q^k p$, за $k \ge 0$.
    \item[б)] $g_X(s) = \frac{p}{1-qs}$
    \item[в)] $\mathbb{E}X = \frac{q}{p}$, $\mathbb{D}X = \frac{q}{p^2}$
\end{enumerate}
\end{theorem}
\begin{proof}
а) Събитието $\{X=k\}$ означава, че първите $k$ експеримента са неуспешни, а $k+1$-вият е успешен. Поради независимостта:
\[ \mathbb{P}(X=k) = \mathbb{P}(X_1=0) \dots \mathbb{P}(X_k=0) \mathbb{P}(X_{k+1}=1) = q^k p. \]

б) За пораждащата функция имаме:
\[ g_X(s) = \sum_{n=0}^{\infty} s^n \mathbb{P}(X=n) = \sum_{n=0}^{\infty} s^n q^n p = p \sum_{n=0}^{\infty} (qs)^n = \frac{p}{1-qs}. \]

в) Чрез диференциране на $g_X(s)$:
\[ g_X'(s) = \frac{pq}{(1-qs)^2} \implies \mathbb{E}X = g_X'(1) = \frac{pq}{(1-q)^2} = \frac{q}{p}. \]
Аналогично, чрез втората производна се получава дисперсията $\mathbb{D}X = \frac{q}{p^2}$.
\end{proof}

\subsection{Свойство безпаметност (Memoryless property)}
Геометричното разпределение е единственото дискретно разпределение, което "няма памет". Тоест, фактът, че вече са настъпили определен брой неуспехи, не променя вероятността за бъдещите експерименти.

\begin{theorem}[Безпаметност]
Ако $X \sim \text{Ge}(p)$, то за $\forall m \ge 0$ и $\forall k \ge 0$:
\[ \mathbb{P}(X \ge m+k \mid X \ge k) = \mathbb{P}(X \ge m). \]
\end{theorem}
\begin{proof}
\[ \mathbb{P}(X \ge m+k \mid X \ge k) = \frac{\mathbb{P}(X \ge m+k \text{ и } X \ge k)}{\mathbb{P}(X \ge k)} = \frac{\mathbb{P}(X \ge m+k)}{\mathbb{P}(X \ge k)}. \]
Тъй като $\mathbb{P}(X \ge k) = q^k$, получаваме $\frac{q^{m+k}}{q^k} = q^m = \mathbb{P}(X \ge m)$.
\end{proof}

\section{Отрицателно биномно разпределение}

Отрицателното биномно разпределение обобщава геометричното. То брои броя на неуспехите до настъпването на $r$-тия успех в схема на Бернули. Означаваме го с $X \sim \text{NB}(r, p)$, където $r \ge 1$ е цяло число.
Забележете, че $\text{NB}(1, p) = \text{Ge}(p)$.
Случайната величина може да се представи като сума от независими геометрични величини: $X = \sum_{j=1}^r Y_j$, където $Y_j \sim \text{Ge}(p)$.

\begin{theorem}
Нека $X \sim \text{NB}(r, p)$. Тогава:
\begin{enumerate}
    \item[а)] $g_X(s) = \left( \frac{p}{1-qs} \right)^r$
    \item[б)] $\mathbb{E}X = r\frac{q}{p}$, $\mathbb{D}X = r\frac{q}{p^2}$
    \item[в)] $\mathbb{P}(X=k) = \binom{r+k-1}{k} p^r q^k$, за $k \ge 0$.
\end{enumerate}
\end{theorem}
\begin{proof}
Свойствата следват директно от факта, че $X$ е сума на $r$ независими величини, разпределени геометрично.
Например, $\mathbb{E}X = \sum_{j=1}^r \mathbb{E}Y_j = r\frac{q}{p}$. 
Вероятностите могат да се получат чисто аналитично чрез $k$-тата производна на $g_X(s)$ в нулата.
\end{proof}

\section{Поасоново разпределение}

Поасоновото разпределение често се използва за моделиране на събития, настъпващи рядко в непрекъснато време или пространство (напр. брой катастрофи, брой клиенти за единица време).

\begin{definition}
Случайната величина $X$ е разпределена по Поасон с параметър $\lambda > 0$, $X \sim \text{Poi}(\lambda)$, ако вероятностите ѝ са дадени от:
\[ \mathbb{P}(X=k) = \frac{\lambda^k}{k!} e^{-\lambda}, \quad \text{за } k \ge 0. \]
\end{definition}
Коректността на това разпределение следва от развитието на експоненциалната функция в ред на Тейлър:
\[ \sum_{k=0}^{\infty} \mathbb{P}(X=k) = e^{-\lambda} \sum_{k=0}^{\infty} \frac{\lambda^k}{k!} = e^{-\lambda} e^{\lambda} = 1. \]

\begin{theorem}
Нека $X \sim \text{Poi}(\lambda)$. Тогава $g_X(s) = e^{\lambda s - \lambda}$, и $\mathbb{E}X = \mathbb{D}X = \lambda$.
\end{theorem}
\begin{proof}
За пораждащата функция имаме:
\[ g_X(s) = \sum_{n=0}^{\infty} s^n \frac{\lambda^n}{n!} e^{-\lambda} = e^{-\lambda} \sum_{n=0}^{\infty} \frac{(s\lambda)^n}{n!} = e^{-\lambda} e^{\lambda s} = e^{\lambda s - \lambda}. \]
Диференцирайки, получаваме $g_X'(s) = \lambda e^{\lambda s - \lambda}$, което при $s=1$ дава $\mathbb{E}X = \lambda$. 
Втората производна е $g_X''(s) = \lambda^2 e^{\lambda s - \lambda}$, откъдето $\mathbb{D}X = \lambda^2 + \lambda - (\lambda)^2 = \lambda$.
\end{proof}

\end{document}